\Conclusion

В~отчёте выполнено описание проекта OneVoice по пунктам чек-листа.

В~результате работы представлены следующие разделы:

\begin{compactlist}
  \item Проблема и целевая аудитория: сформулирована проблема рассогласования и ручного управления цифровым присутствием МСБ на множественных платформах; целевая аудитория~--- субъекты малого и среднего бизнеса (типовой пример~--- кофейня).
  \item Текущее состояние: зафиксирована модель AS-IS~--- ручное управление на каждой платформе, отсутствие эталонного источника данных, разрозненный мониторинг отзывов и сообщений.
  \item Аналоги и конкуренты: выделены классы решений (CRM, SMM, ORM), приведены критерии сравнения (К1--К6) и балльная таблица; показано, что ни один аналог не закрывает задачу консистентности данных.
  \item Предметная область: описано цифровое присутствие, консистентность данных и репутационный контур; исследовательская часть~--- классификация платформ по типу интеграции (API / RPA).
  \item Практическая часть и технологии: указан технологический стек (Go, Next.js, PostgreSQL, MongoDB, NATS, Playwright), протокол A2A, спецификации агентов (VK, Telegram, Яндекс.Бизнес).
  \item Данные: перечислены сущности (профили бизнеса, интеграции, отзывы, задачи), источники данных и требования к~обработке персональных данных (152-ФЗ).
\end{compactlist}

Цель работы~--- разработка программного комплекса на основе мультиагентной архитектуры с~гибридной моделью интеграции для управления цифровым присутствием МСБ~--- отражена в~постановке задачи и требованиях; перечень задач приведён во введении в~хронологическом порядке. Отчёт готов для использования в~рамках сопровождения ВКР.
